% GENERATED FILE. Edit canonical lessons, then run npm run generate:books.
% canonical-source-hash: fnv1a64:84e926110f3ae983
% canonical-lessons: TE-C149-naducu, TE-C149-muyu, TE-C149-teruvu, TE-C149-tagu, TE-C149-edcu, TE-R149-first-pass-describing-words-and-verbs, TE-R149-second-pass-verbs

\chapter{To walk, To close, To open, To drink, To cry}
\label{ch:te-to-walk-to-close-to-open-to-drink-to-cry}
\hlchaptermodality{149}

\begin{chapteropening}
\textbf{By the end of this chapter:} \emph{I can say to walk, to close, to open, to drink and to cry.}

Complete the last lesson of chapter 149: I can say to walk, to close, to open, to drink and to cry.
\end{chapteropening}

\section[\texorpdfstring{\te{నడుచు} (naḍucu)}{నడుచు (naḍucu)}]{\te{నడుచు} (naḍucu) — to walk}

\label{lesson:TE-C149-naducu}



\begin{quote}
Before the new one: say the Telugu for to sing, then the Telugu for to sell.
\end{quote}

\subsection*{You'll want to know: \te{నడుచు}}
\textbf{\te{నడుచు}} — \emph{naḍucu} — \textquotedblleft{}to walk\textquotedblright{}.

\begin{grammarlens}[title={What you've built}]
One more word for this chapter.
\end{grammarlens}

\subsection*{Your turn}
\begin{itemize}
\raggedright
  \item \emph{Say it:} \emph{naḍucu}
  \item \emph{Say it:} \emph{naḍucu}, once more
  \item \emph{Say it:} say \emph{naḍucu}
  \item \emph{Recall:} say the Telugu for to sing, then the Telugu for to sell, then say \emph{naḍucu} again
\end{itemize}

\begin{tcolorbox}[breakable,colback=teal!4,colframe=teal!35!black,title={Before you move on}]
What does \te{నడుచు} mean? (\textquotedblleft{}To walk\textquotedblright{}.) Say it once more.
\end{tcolorbox}
\section[\texorpdfstring{\te{మూయు} (mūyu)}{మూయు (mūyu)}]{\te{మూయు} (mūyu) — to close}

\label{lesson:TE-C149-muyu}



\begin{quote}
Before the new one: say the Telugu for to sell, then the Telugu for to walk.
\end{quote}

\subsection*{You'll want to know: \te{మూయు}}
\textbf{\te{మూయు}} — \emph{mūyu} — \textquotedblleft{}to close\textquotedblright{}.

\begin{grammarlens}[title={What you've built}]
One more word for this chapter.
\end{grammarlens}

\subsection*{Your turn}
\begin{itemize}
\raggedright
  \item \emph{Say it:} \emph{mūyu}
  \item \emph{Say it:} \emph{mūyu}, once more
  \item \emph{Say it:} say \emph{mūyu}
  \item \emph{Recall:} say the Telugu for to sell, then the Telugu for to walk, then say \emph{mūyu} again
\end{itemize}

\begin{tcolorbox}[breakable,colback=teal!4,colframe=teal!35!black,title={Before you move on}]
What does \te{మూయు} mean? (\textquotedblleft{}To close\textquotedblright{}.) Say it once more.
\end{tcolorbox}
\section[\texorpdfstring{\te{తెరువు} (teruvu)}{తెరువు (teruvu)}]{\te{తెరువు} (teruvu) — to open}

\label{lesson:TE-C149-teruvu}



\begin{quote}
Before the new one: say the Telugu for to walk, then the Telugu for to close.
\end{quote}

\subsection*{You'll want to know: \te{తెరువు}}
\textbf{\te{తెరువు}} — \emph{teruvu} — \textquotedblleft{}to open\textquotedblright{}.

\begin{grammarlens}[title={What you've built}]
One more word for this chapter.
\end{grammarlens}

\subsection*{Your turn}
\begin{itemize}
\raggedright
  \item \emph{Say it:} \emph{teruvu}
  \item \emph{Say it:} \emph{teruvu}, once more
  \item \emph{Say it:} say \emph{teruvu}
  \item \emph{Recall:} say the Telugu for to walk, then the Telugu for to close, then say \emph{teruvu} again
\end{itemize}

\begin{tcolorbox}[breakable,colback=teal!4,colframe=teal!35!black,title={Before you move on}]
What does \te{తెరువు} mean? (\textquotedblleft{}To open\textquotedblright{}.) Say it once more.
\end{tcolorbox}
\section[\texorpdfstring{\te{తాగు} (tāgu)}{తాగు (tāgu)}]{\te{తాగు} (tāgu) — to drink}

\label{lesson:TE-C149-tagu}



\begin{quote}
Before the new one: say the Telugu for to close, then the Telugu for to open.
\end{quote}

\subsection*{You'll want to know: \te{తాగు}}
\textbf{\te{తాగు}} — \emph{tāgu} — \textquotedblleft{}to drink\textquotedblright{}.

\begin{grammarlens}[title={What you've built}]
One more word for this chapter.
\end{grammarlens}

\subsection*{Your turn}
\begin{itemize}
\raggedright
  \item \emph{Say it:} \emph{tāgu}
  \item \emph{Say it:} \emph{tāgu}, once more
  \item \emph{Say it:} say \emph{tāgu}
  \item \emph{Recall:} say the Telugu for to close, then the Telugu for to open, then say \emph{tāgu} again
\end{itemize}

\begin{tcolorbox}[breakable,colback=teal!4,colframe=teal!35!black,title={Before you move on}]
What does \te{తాగు} mean? (\textquotedblleft{}To drink\textquotedblright{}.) Say it once more.
\end{tcolorbox}
\section[\texorpdfstring{\te{ఏడ్చు} (ēḍcu)}{ఏడ్చు (ēḍcu)}]{\te{ఏడ్చు} (ēḍcu) — to cry}

\label{lesson:TE-C149-edcu}



\begin{quote}
Before the new one: say the Telugu for to open, then the Telugu for to drink.
\end{quote}

\subsection*{You'll want to know: \te{ఏడ్చు}}
\textbf{\te{ఏడ్చు}} — \emph{ēḍcu} — \textquotedblleft{}to cry\textquotedblright{}.

\begin{grammarlens}[title={What you've built}]
One more word for this chapter.
\end{grammarlens}

\subsection*{Your turn}
\begin{itemize}
\raggedright
  \item \emph{Say it:} \emph{ēḍcu}
  \item \emph{Say it:} \emph{ēḍcu}, once more
  \item \emph{Say it:} say \emph{ēḍcu}
  \item \emph{Recall:} say the Telugu for to open, then the Telugu for to drink, then say \emph{ēḍcu} again
\end{itemize}

\begin{tcolorbox}[breakable,colback=teal!4,colframe=teal!35!black,title={Before you move on}]
What does \te{ఏడ్చు} mean? (\textquotedblleft{}To cry\textquotedblright{}.) Say it once more.
\end{tcolorbox}
\section[(dialogue)]{Describing words and verbs — a first pass}

\label{lesson:TE-R149-first-pass-describing-words-and-verbs}



\begin{quote}
Say the words for to drink and to cry.
\end{quote}

\subsection*{Your turn}
Say these aloud:

\begin{itemize}
\raggedright
  \item safe, a lie; untrue, grey; ash, golden, spoiled
  \item sweet, bitter-tasting, spicy, courage; brave, lazy
  \item hard-working, clever, terrible, black, white
  \item red, to fly, to run, to sleep, to get up
  \item to cook, to cut, to break, to send, to look for
\end{itemize}

\begin{tcolorbox}[breakable,colback=teal!4,colframe=teal!35!black,title={Before you move on}]
Safe? (\textbf{\te{సురక్షితం}}, \emph{surakṣitaṁ}.) A lie; untrue? (\textbf{\te{అబద్ధం}}, \emph{abaddhaṁ}.) Grey; ash? (\textbf{\te{బూడిద}}, \emph{būḍida}.) Golden? (\textbf{\te{బంగారు}}, \emph{baṅgāru}.) Spoiled? (\textbf{\te{పాడైన}}, \emph{pāḍaina}.) Sweet? (\textbf{\te{తీయని}}, \emph{tīyani}.) Bitter-tasting? (\textbf{\te{చేదైన}}, \emph{cēdaina}.) Spicy? (\textbf{\te{కారమైన}}, \emph{kāramaina}.) Courage; brave? (\textbf{\te{ధైర్యం}}, \emph{dhairyaṁ}.) Lazy? (\textbf{\te{సోమరి}}, \emph{sōmari}.) Hard-working? (\textbf{\te{కష్టపడే}}, \emph{kaṣṭapaḍē}.) Clever? (\textbf{\te{తెలివైన}}, \emph{telivaina}.) Terrible? (\textbf{\te{భయంకరమైన}}, \emph{bhayaṅkaramaina}.) Black? (\textbf{\te{నల్లని}}, \emph{nallani}.) White? (\textbf{\te{తెల్లని}}, \emph{tellani}.) Red? (\textbf{\te{ఎర్రటి}}, \emph{erraṭi}.) To fly? (\textbf{\te{ఎగురు}}, \emph{egiru}.) To run? (\textbf{\te{పరిగెత్తు}}, \emph{parigettu}.) To sleep? (\textbf{\te{నిద్రపో}}, \emph{nidrapō}.) To get up? (\textbf{\te{లేచు}}, \emph{lēcu}.) To cook? (\textbf{\te{వండు}}, \emph{vaṇḍu}.) To cut? (\textbf{\te{కోయు}}, \emph{kōyu}.) To break? (\textbf{\te{పగలగొట్టు}}, \emph{pagalagoṭṭu}.) To send? (\textbf{\te{పంపు}}, \emph{pampu}.) To look for? (\textbf{\te{వెతుకు}}, \emph{vetuku}.)
\end{tcolorbox}
\section[(dialogue)]{Verbs — a second pass}

\label{lesson:TE-R149-second-pass-verbs}



\begin{quote}
Say the words for to learn and to teach.
\end{quote}

\subsection*{Your turn}
Say these aloud:

\begin{itemize}
\raggedright
  \item to learn, to teach, to forget, to jump, to get down
  \item to reach, to put on, to call, to count, to wash
  \item to wipe, to tie, to sew, to win, to hide
  \item to protect, to burn, to put out, to sing, to sell
  \item to walk, to close, to open, to drink, to cry
\end{itemize}

\begin{tcolorbox}[breakable,colback=teal!4,colframe=teal!35!black,title={Before you move on}]
To learn? (\textbf{\te{నేర్చుకో}}, \emph{nērcukō}.) To teach? (\textbf{\te{నేర్పు}}, \emph{nērpu}.) To forget? (\textbf{\te{మరచిపో}}, \emph{maracipō}.) To jump? (\textbf{\te{దూకు}}, \emph{dūku}.) To get down? (\textbf{\te{దిగు}}, \emph{digu}.) To reach? (\textbf{\te{చేరు}}, \emph{cēru}.) To put on? (\textbf{\te{వేసుకో}}, \emph{vēsukō}.) To call? (\textbf{\te{పిలుచు}}, \emph{pilucu}.) To count? (\textbf{\te{లెక్కపెట్టు}}, \emph{lekkapeṭṭu}.) To wash? (\textbf{\te{కడుగు}}, \emph{kaḍugu}.) To wipe? (\textbf{\te{తుడుచు}}, \emph{tuḍucu}.) To tie? (\textbf{\te{కట్టు}}, \emph{kaṭṭu}.) To sew? (\textbf{\te{కుట్టు}}, \emph{kuṭṭu}.) To win? (\textbf{\te{గెలుచు}}, \emph{gelucu}.) To hide? (\textbf{\te{దాచు}}, \emph{dācu}.) To protect? (\textbf{\te{కాపాడు}}, \emph{kāpāḍu}.) To burn? (\textbf{\te{కాల్చు}}, \emph{kālcu}.) To put out? (\textbf{\te{ఆర్పు}}, \emph{ārpu}.) To sing? (\textbf{\te{పాడు}}, \emph{pāḍu}.) To sell? (\textbf{\te{అమ్ము}}, \emph{ammu}.) To walk? (\textbf{\te{నడుచు}}, \emph{naḍucu}.) To close? (\textbf{\te{మూయు}}, \emph{mūyu}.) To open? (\textbf{\te{తెరువు}}, \emph{teruvu}.) To drink? (\textbf{\te{తాగు}}, \emph{tāgu}.) To cry? (\textbf{\te{ఏడ్చు}}, \emph{ēḍcu}.)
\end{tcolorbox}
